\begin{table}[!htbp]
\centering
\caption{A13 follow-up variance and effect-estimation diagnostics.}
\label{tab:supp-focused-A13}
\begin{threeparttable}
\begin{singlespace}
\small
\setlength{\tabcolsep}{3.5pt}
\renewcommand{\arraystretch}{1.1}
\begin{tabular}{@{}lrr@{}}
\toprule
Quantity & Unadjusted & Adjusted\\
\midrule
$\overline{\widehat V_{\rm S}}/\widehat{\Var}(T_n)$ & 0.9995 & 1.0019 \\
$\overline{V^*}/\widehat{\Var}(T_n)$ & 1.0004 & 1.0027 \\
$\overline{\widehat V_{\rm R}^{\rm raw}}/\overline{\widehat V_{\rm S}}$ & 1.0062 & 1.0065 \\
Median $\widehat\lambda$ & 1.0024 & 1.0025 \\
95th percentile of $\widehat\lambda$ & 1.0119 & 1.0124 \\
CRT using true $\bm d=0$ (\%) & 4.969 & 4.937 \\
Exploratory noise-subtracted CRT (\%) & 5.069 & 5.074 \\
Quadratic noise / mean $\widehat V_{\rm S}$ (\%) & 0.624 & 0.652 \\
Estimated subtraction / mean $\widehat V_{\rm S}$ (\%) & 0.473 & 0.494 \\
Singleton expected noise / mean $\widehat V_{\rm S}$ (\%) & 0.147 & 0.154 \\
\bottomrule
\end{tabular}
\begin{tablenotes}[flushleft]\footnotesize
\item[] All entries use 100,000 trials and reconstructed covariance. $V^*$ is the sample variance of the regenerated reference statistics. The true-effect and noise-subtracted procedures are diagnostics, not replacements for the principal CRT. The variance ratios are ratios of means; the quantiles describe the trial-specific scale factor.
\end{tablenotes}
\end{singlespace}
\end{threeparttable}
\end{table}
